\section*{Capítulo \ref{ch:b3:electronic-spectroscopy} --- Espectroscopia eletrônica e fotofísica}

\begin{solution}{exo:b3:electronic-spectroscopy:1}
(a) Proibida pela regra de spin ($\Delta S = 1$). (b) Proibida pela \omterm{thm:b3:electronic-spectroscopy:laporte}{regra de Laporte}
($g \to g$); vista fracamente graças ao acoplamento vibrônico. (c) Permitida: $B_{1u}$ é a
representação de $z$ em $D_{2h}$. (d) Proibida por simetria: $A_2$ não é nem $x$, nem
$y$, nem $z$ em $C_{2v}$; a banda é fraca.
\end{solution}

\begin{solution}{exo:b3:electronic-spectroscopy:2}
$10^7/349 - 10^7/450 = 28653 - 22222 = \qty{6431}{cm^{-1}}$, \qty{0.80}{eV}.
\end{solution}

\begin{solution}{exo:b3:electronic-spectroscopy:3}
$A = 5700 \times 1 \times 1.0\times10^{-5} = 0.057$; fração absorvida $1 - 10^{-0.057} =
0.12$.
\end{solution}

\begin{solution}{exo:b3:electronic-spectroscopy:4}
$k_r = \Phi_F/\tau = \qty{1.5e8}{s^{-1}}$, $\tau_r = \qty{6.7}{ns}$; $k_{nr} = (1 - \Phi_F)/\tau
= \qty{1.0e8}{s^{-1}}$.
\end{solution}

\begin{solution}{exo:b3:electronic-spectroscopy:5}
$\int\varepsilon\,\dd\tilde\nu = 1.0645 \times 5700 \times 5000 = \num{3.03e7}$, logo $f \approx
4.32\times10^{-9} \times 3.03\times10^7 = 0.13$: uma transição permitida de intensidade
moderada.
\end{solution}

\begin{solution}{exo:b3:electronic-spectroscopy:6}
O fator de Poisson $\eu^{-S}S^v/v!$ é máximo em $v = \lfloor S\rfloor$. $S = 0.3$: a
própria raia 0--0. $S = 1.5$: $v' = 1$, 1.5 vez a raia 0--0. $S = 5$: $v' = 4$
e 5, igualmente, 26 vezes a raia 0--0.
\end{solution}

\begin{solution}{exo:b3:electronic-spectroscopy:7}
$\sum k = \qty{4.2e8}{s^{-1}}$: $\tau = \qty{2.4}{ns}$, $\Phi_F = 0.24$, $\Phi_{\mathrm{ISC}} = 0.71$
(e $\Phi_{\mathrm{IC}} = 0.05$).
\end{solution}

\begin{solution}{exo:b3:electronic-spectroscopy:8}
$\Phi = 0.546 \times 0.46 \times (0.050/0.040) \times (1.4266/1.333)^2 = 0.36$.
\end{solution}

\begin{solution}{exo:b3:electronic-spectroscopy:9}
$\tau_0/\tau = 1.429$ e 1.860: linear em $[Q]$, logo o próprio tempo de vida é
encurtado: \omterm{def:b3:electronic-spectroscopy:quencher}{supressão dinâmica}. $K_{SV} = 43$ L/mol (os dois pontos), $k_q = 43/8.0
\times 10^{-9} = \qty{5.4e9}{L.mol^{-1}.s^{-1}}$.
\end{solution}

\begin{solution}{exo:b3:electronic-spectroscopy:10}
\omterm{def:b3:electronic-spectroscopy:quencher}{Supressão estática}: as moléculas que ainda emitem vivem tanto quanto antes; metade
delas está ligada, no estado fundamental, em um complexo não fluorescente. Com uma
constante de associação $K_a$, a fração livre é $1/(1 + K_a[Q])$ e $I_0/I = 1 +
K_a[Q]$, a mesma forma que Stern--Volmer, mas com $\tau$ inalterado.
\end{solution}

\begin{solution}{exo:b3:electronic-spectroscopy:11}
$\tau_r = \tau/\Phi_F = \qty{14}{ns}$; $\Phi_T = 1 - 0.36 = 0.64$. O tempo de vida medido é
encurtado pelo \omterm{def:b3:electronic-spectroscopy:radiationless}{cruzamento intersistemas} concorrente; $\tau_r$ é o que a emissão
sozinha daria.
\end{solution}

\begin{solution}{exo:b3:electronic-spectroscopy:12}
$\frac12\langle\alpha\beta - \beta\alpha|\alpha\beta + \beta\alpha\rangle = \frac12(1 + 0 - 0 - 1) = 0$, usando
$\langle\alpha|\alpha\rangle = \langle\beta|\beta\rangle = 1$, $\langle\alpha|\beta\rangle = 0$ para cada elétron.
O \omterm{def:b3:quantum-model-systems:operator}{operador} dipolo não toca o spin, logo o momento de transição contém esse fator
nulo: $T_1 \leftarrow S_0$ não tem intensidade, a menos que o \omterm{def:b3:many-electron-atoms:spin-orbit}{acoplamento spin--órbita}
misture os estados.
\end{solution}

\begin{solution}{pb:b3:electronic-spectroscopy:1}
\textbf{1.} $A = 5700 \times 1 \times 1.0\times10^{-4} = 0.57$.
\textbf{2.} $1 - 10^{-0.57} = 0.73$.
\textbf{3.} $\varepsilon$ de alguns milhares: uma transição $\pi\to\pi^*$ permitida, de
intensidade moderada.
\textbf{4.} $1239.8/349 = \qty{3.55}{eV}$.
\textbf{5.} Ela absorve apenas no ultravioleta; a luz visível passa.
\textbf{6.} Para manter a luz absorvida proporcional à concentração e evitar a
reabsorção da luz emitida (efeitos de filtro interno).
\textbf{7.} $S_0 \to S_1$ (absorção, depois \omterm{def:b3:electronic-spectroscopy:radiationless}{relaxação vibracional}); a partir de $S_1$:
\omterm{def:b3:electronic-spectroscopy:luminescence}{fluorescência}, \omterm{def:b3:electronic-spectroscopy:radiationless}{conversão interna}, \omterm{def:b3:electronic-spectroscopy:radiationless}{cruzamento intersistemas} para $T_1$.
\textbf{8.} $k_r = 0.546/19 \times 10^{-9} = \qty{2.87e7}{s^{-1}}$; $\tau_r = \qty{35}{ns}$.
\textbf{9.} $(1 - 0.546)/\tau_0 = \qty{2.39e7}{s^{-1}}$.
\textbf{10.} $28653 - 22222 = \qty{6431}{cm^{-1}}$.
\textbf{11.} A emissão parte de $S_1$ relaxado e termina em níveis vibracionalmente
excitados de $S_0$, depois que a molécula e o solvente relaxaram: o fóton tem menos
energia, aqui o bastante para cair no visível.
\textbf{12.} $\Phi_F \times (349/450) = 0.42$: 42\,\% da energia absorvida.
\textbf{13.} Os pontos sobem 0.594 a cada \qty{0.010}{mol/L}: uma reta que passa por 1.
\textbf{14.} Mínimos quadrados com os cinco pontos: inclinação $K_{SV} = \qty{59.4}{L/mol}$,
intercepto 1.000.
\textbf{15.} $k_q = K_{SV}/\tau_0 = \qty{3.1e9}{L.mol^{-1}.s^{-1}}$.
\textbf{16.} $\tau = \tau_0/(1 + 59.4 \times 0.040) = \qty{5.6}{ns}$.
\textbf{17.} Medindo os tempos de vida: na \omterm{def:b3:electronic-spectroscopy:quencher}{supressão dinâmica}, $\tau_0/\tau$ cai sobre a
mesma reta que $I_0/I$.
\textbf{18.} A água tônica não contém cloreto que suprima a \omterm{def:b3:electronic-spectroscopy:luminescence}{fluorescência}, e é ácida o
bastante para manter o quinino protonado, sua forma fluorescente.
\textbf{19.} $k_d = 8 \times 8.314 \times 298/(3 \times 0.890\times10^{-3}) =
\qty{7.4e6}{m^3.mol^{-1}.s^{-1}} = \qty{7.4e9}{L.mol^{-1}.s^{-1}}$.
\textbf{20.} $k_q$ é um pouco menos da metade de $k_d$.
\textbf{21.} Cerca de 42\,\%.
\textbf{22.} Mais lenta: $k_d \propto 1/\eta$, e $k_q$ só pode cair junto.
\textbf{23.} \textbf{$k_q/k_d = 0.42$}: o cloreto suprime a \omterm{def:b3:electronic-spectroscopy:luminescence}{fluorescência} do quinino a quase
metade da velocidade dos encontros controlados pela difusão.
\end{solution}
